

\prSect{INTRODUCTION}

\pfl{}

We begin our updated critique of the American family court system by
acknowledging that, notwithstanding the fact that mothers are often
placed at risk of incarceration, fines, and losing all contact with their
children in retaliation for requesting the family courts to shield their
children from abuse committed by the other parent, there are still mea-
sured successes in this uphill battle to reform the family courts. Such
successes give us reason for hope. By elucidating the problems preva-
lently found in the family court system, we see evidence of steady efforts
to effect change. For those who have become aware of the troubling pat-
terns of family court dysfunction, it becomes increasingly hard for them
to accept the dismal status quo. \pfl{}

Most notably, a cadre of reformers have set the wheels in motion for
course-correcting a system that has careened out of control in those—all
too frequent—instances in which family court decisions are putatively
and provably contraindicative to the safety and well-being of abused
children. Taking a balanced approach, we acknowledge the fact that
there are many good judges who zealously protect children from harm.
However, there are many who do not. And too many children in custody
matters are thus falling through the cracks. \pfl{}

Among those who have had a significant effect on improving the
family court and child welfare system is Dr. Jeffrey L. Edleson, whose
noteworthy contributions will be discussed at length in the addendum to
\notePage{xxxiii}
the third section of this book. Specifically, Dr. Edleson, in collaboration
with his colleague, Susan Schechter, authored a comprehensive guide
for court professionals and child advocates in protecting the welfare and
safety of children. Their seminal book, Effective Intervention in Domestic
Violence and Child Maltreatment: Guidelines for Policy and Practice 1 —
better known as the “Greenbook” (after the color of its cover)—whose
focus is “to present leaders of communities and institutions with a
context-setting tool to develop public policy aimed at keeping families
safe and stable,” 2 has effectively spawned an impressive body of literature
based on its comprehensive and cogent recommendations for guiding
professionals in working with battered women and abused children. \pfl{}

Although first written in 1999, a concrete example of the effects of
the Greenbook in improving child safety can be found twenty years later
in a 2019 journal article. In “Advocacy Perspectives on the Greenbook,”
Lucy Salcido Carter explained that “[w]hen domestic violence is an issue
in the case, it is important not to have [issued] an automatic finding
of failure to protect against the non-offending parent [the mother], but
instead to recognize that abuse is happening to the adult victim and that
the perpetrator is exposing the child to the domestic violence.” 3 \pfl{}

Similar to the Greenbook’s introduction of helpful policies and prac-
tices in protecting children, on the clinical front some noteworthy
advances have been made that deserve pointing out. Most notable, emi-
nent neuropsychologist and family counseling expert Dr. Robert (Bob)
Geffner spearheads a unique continuing education training program
for clinicians that introduces advanced research and clinical methods
for assessing and treating intrafamilial abuse and trauma. Founder of
the Institute on Violence, Abuse and Trauma, (IVAT), under Geffner’s
stewardship thousands of therapists have received training in abuse
and trauma intervention since its inception in 1984. With a focus on
educating the mental health profession to the traumatic effects on chil-
dren judicially ordered into the custody of the sexually abusive parent,
Geffner has paved the way for children and adult survivors to receive
the trauma intervention therapy to heal from unremitting exposure to
ongoing assault. His work has resonated throughout the mental health
profession and has established a new standard. \pfl{}

\notePage{xxxiv}
All in all, dedicated researchers and clinicians have taught us that it re-
quires patience and forbearance to rehabilitate a system that has become
woefully tainted by inimical prejudices and bad practices. We hope this
painstakingly candid critique of the family court system will motivate,
inspire, and galvanize change in how judges approach the complexity of
custody decisions when there are credible charges of child sexual abuse
leveled against the other parent—so that tendentious mental health the-
ories interlaced with gender, racial, ethnic, and social class biases do not
inform ill-fated family court rulings. \pfl{}


We must also bear in mind that it is imperative to enlarge our army
of family court reformers, lest our successes remain limited and iso-
lated at best. While acknowledging the small but measurable successes
in enlightening judges and others to the pitfalls of system biases that
severely compromise a child’s welfare and safety, we don’t deny the ex-
tant horrors either. In case after case, the evidence shows the system
maniacally works against mothers and throws abused children into a
culvert of despair. In fact, a common phrase heard today by mothers
caught in family court litigation is that they feel beleaguered by a “trial
by ambush.” They are referring to the unremitting onslaught of legal
attacks—cacophonous winds of frivolous legal filings—which all too of-
ten entail the Attorney for the Child (AFC) and Child Protective Services
(CPS) who, in concert with the abusive parent, target the mother without
pause. Sadly, these horrific family court cases are all too common—they
are neither anomalies nor aberrations. In fact, the intensity of the per-
secution of mothers in the family courts and child welfare system is
seen even more so among Black, Indigenous, People of Color (BIPOC)
populations, who have the added burden of contending with the race
factor in addition to the systemic biases and failures of the family court
system. \pfl{}


For these reasons, we offer a second edition of From Madness to
Mutiny in the hope that it will stimulate purposeful and productive dis-
course in both the classrooms that breed the next generation of lawyers,
judges, social workers, and guardians ad litem, and in the clinical set-
ting wherein everyday matters of domestic violence and child abuse
become evident. Sociologists, psychologists, feminists, policy makers,
\notePage{xxxv}
and others must be enlightened to the scourge of gender bias that has
entrapped far too many mothers in the family courts. Those mothers suf-
fer the anguish of seeing their children harmed; many are subsequently
expunged from their children’s lives or restricted to minimal contact
with them. We must extirpate this malady at its roots. This was the pur-
pose of the first edition of our book, and likewise the warrant for the
second. \pfl{}

\prPart{FAMILY COURT CUSTODY CRISIS OF EPIC PROPORTION}

\prHead{United Nations Addresses Family Courts’ Pervasive Gender-Bias
and the “Double” Victimization of Women and Children as a Global
Human Rights Crisis}


In the fall of 2022, the United Nations Human Rights Office of the
High Commissioner issued a “Call for Inputs—Custody Cases, Violence
against Women and Violence against Children,” which was a broad alert
to academic institutions, international organizations, and governmen-
tal agencies to supply data on the gravity of abuses perpetrated against
women and children caught in the family courts worldwide. \pfl{}


Pointing to family courts’ tendency to dismiss evidence of domestic
violence and child abuse in custody cases, the Call for Inputs asserted
that many women who are victims of violence and whose children have
been abused by the father will “face double victimization as they are
punished for alleging abuse,” and by “losing custody” of their children
“or at times being imprisoned.” Identifying what is at the crux of these
disastrous custody decisions, the UN Office of the High Commissioner
pointed to family courts’ widespread reliance upon concepts such as
“parental alienation” (in which a parent, usually the mother, is blamed,
albeit unfairly, for any hostility shown by an abused child toward its
father, whom the child had already identified as their abuser). They
pointed out that though “these concepts lack a universal clinical or sci-
entific definition,” they are widely accepted in the family courts as the
justification for stripping a mother of custody of her child. \pfl{}

\notePage{xxxvi}
Driving this major UN Human Rights project was a need “to exam-
ine the ways in which family courts in different world regions refer
to parental alienation, or similar concepts, in custody cases and how
this may lead to double victimization of victims of domestic violence or
abuse” and “to document the many ways in which family courts ignore
the history and existence of domestic . . . violence and [child] abuse in
the context of custody cases” while studying “their grave consequences
[to] mothers and their children.” \pfl{}


Furthermore, the Human Rights Office of the High Commissioner
points out that one of the goals of their international study is to learn
more about “any intersecting elements such as age, sex, gender, race,
ethnicity, legal residence, religious or political belief ” that result in poor
case outcome for women and children. They stated the overall purpose
of their study is to “draw attention to the scale and manifestation” of the
double victimization of women and children when courts across “many
countries, spanning all regions of the world” discredit mothers’ pleas
to protect their children from abuse and to “offer recommendations for
States and other stakeholders to address the situation.” 4 \pfl{}


The UN study offered hope to women’s groups across the globe. In
the United States, The Women’s Coalition—a well-organized group of
mothers and advocates devoted to collectively raising awareness about
the custody crisis—was exultant about the project. The group’s founder,
San Diego-based Cindy Dumas, a PhD candidate in psychology whose
career was aborted when she became a victim of the family court and
was forced to go on the run with her children, posted the following on
her website:

\prQuote{This is the first time a mechanism of the UN HRC has priori-
tized this issue and initiated an inquiry. It is a leap forward in our
efforts to get it officially recognized as a human rights violation.
The Women’s Coalition has previously submitted three petitions,
which have likely contributed to them initiating this inquiry.}5

\noindent{}Shortly after the UN’s announcement of their “Call for Inputs,” a ques-
tionnaire was sent out to protective mothers by The Women’s Coalition.
\notePage{xxxvii}
Within the first week or so, Dumas would receive nearly four hundred
completed questionnaires from mothers who lost their children to a
unsympathetic family court system, many of whom lost visitation privi-
leges with their children and, sadly, remain permanently estranged from
them. Her questionnaire covered all the component parts of the family
court system that contribute to the reprehensible rulings against moth-
ers. But even more so, it captured the deleterious effects on children
forced to live with an abusive father where they have become so in-
exorably brainwashed against their mother that reconciliation becomes
impossible. \pfl{}


By early summer 2023, the UN’s Human Rights Office of the High
Commissioner would produce a conclusive report on its findings. In the
press release accompanying their report, they used an alarming heading:
“Urgent reforms needed to protect woman and children from violence
in custody battles: UN expert.” They began with a condemnation of the
family courts: “Deeply embedded gender bias that pervades family court
systems across the globe is placing women and children in situations of
immense suffering and violence, a UN expert said today.” 6 \pfl{}


UN Special Rapporteur, Reem Alsalem, who conducted the study,
went on to state “the tendency of family courts to dismiss the history of
domestic violence and abuse in custody cases, especially where mothers
and/or children have brought forward credible allegations of domestic
abuse, including coercive control, physical or sexual abuse is unaccept-
able.” Perhaps most telling was Alsalem’s sounding the alarm about the
effects of a failed family court system on BIPOC mothers: “The report
underscores that minority women face additional barriers when be-
ing charged of [sic] using ‘parental alienation’ in part due to increased
barriers in accessing justice as well as negative stereotypes.” 7 \pfl{}

\prHead{Inside the Torture Chambers of the Family Courts: Forced
Weaning of Young Infants to Accommodate Overnight Visits
to the Abusive Father}

In recent years, family courts have been removing very young children,
including nursing infants, from mothers with much greater frequency
\notePage{xxxviii}
than they had in the past. In truth, we did study cases where infants were
removed from protective mothers when we conducted research for the
first edition. However, today the family courts are found to strip mothers
of toddlers and infants much more often than they did twenty years ago.
In fact, in preparation for transferring custody of an infant to the father,
the court sets a plan into action by first ordering the infant to go on
extensive overnight visitations (three to four times a week) and weekend
visits with the father as preparation for the transfer of custody to the
father. To comply, the nursing mother must use a pump to express her
milk so that their infant can be sent over to the father with the mother’s
milk supplied. \pfl{}


Such an arrangement is certainly not without hazards. Mothers, for
example, in all kinds of inclement weather, can be seen bundling up their
nursing infants, accompanied by bottles of pumped/expressed breast
milk, as they deliver their babies to the sexually abusive father on the
average of three to four nights a week and on weekends as well. Given
the reality that many nursing infants often cannot adjust that easily to
expressed milk (or formula), they consequently go hungry because they
are used to taking nourishment from their mother’s breast and not from
a bottle. And, sometimes, when their nursing infants are returned to
them from an overnight or weekend visit with the father, the mothers
will notice that their baby’s diaper will show streaks of blood. That is
caused by the infant’s dehydration resulting from their refusal to take
expressed milk or formula, both of which are delivered in the form of a
bottle as opposed to the breast. \pfl{}


Mothers caught in this bind will often prematurely wean their infant
from breastmilk altogether so that the baby will get used to bottled milk
when court-ordered to go to the abusive father for overnights and week-
ends. While this prevents dehydration because the baby has now been
weaned (albeit prematurely) and can therefore take milk from a bottle,
such premature weaning is in direct contravention to the health guide-
lines of today. Both the World Health Organization (WHO) and the
American Academy of Pediatrics recommend breast feeding of an in-
fant for at least one year, and ideally for two to two-and-a-half years.
Breastfeeding inures to the benefit of the baby both emotionally and
physically by strengthening its immune system to fight off disease,
\notePage{xxxix}
improving its growth and development, and fostering a healthy
maternal–child bond. \pfl{}


In the winter of 2023, a Virginia mother made headlines news when
the judge forced her to wean her four-month-old infant so that the in-
fant’s father could have extensive visitation with the child, even though
he had abandoned the mother during her pregnancy. 8 The Washington
Post, which broke the story, set off a firestorm when the fathers’ rights ac-
tivists falsely claimed the mother was using breastfeeding as a “weapon
against visitation.” One of the authors of this book (Amy Neustein)
spoke to the mother at length and learned of her Sisyphean struggle
against a court system that directed asperities at her, freely accusing her
of being “uncooperative.” The facts placed on the record clearly belied
these freewheeling accusations launched at the mother, who did every-
thing in her power to encourage the father’s participation in the life of
the child. In fact, she pleaded with him to allow her to breastfeed her
infant by making arrangements for him to come over to her home for
many hours so that she could nurse the baby while he would visit. The
court found that unacceptable and ordered the mother to express her
milk and send the baby off to be with the father for long stretches of
time. \pfl{}


Concerned about the mistreatment of this lactating mother,
Neustein—together with author Michelle Etlin—wrote a letter to The
Washington Post. The paper published their letter under the heading
“Infants are being penalized in the name of ‘men’s rights.’” Their letter
begins with a sad but true commentary: “This story is all too familiar to
those who study the family court system. For the past four decades, we
have researched how U.S. domestic relations courts . . . have cruelly and
senselessly torn infants and children from their mothers.” 9 \pfl{}


\prHead{Another Instance of Violence against Women: Jailing of Mothers
on Untenable Grounds}

Among the other hardships we have seen in recent years imposed by
courts is the upsurge in jailing of mothers for nonsensical reasons, even
\notePage{xl}
at the cost of their health. A protective mother in Maryland, for example,
was going into insulin shock and was therefore unable to pay attention
to the judge’s questions. The judge became furious and threw her into
jail for the day. A Westchester County, New York, mother was thrown
in jail because she came late for her court hearing, while having driven
across town in bumper-to-bumper traffic. \pfl{}


Mothers, even if not jailed, have in recent years been threatened with
incarceration for talking to their own supportive family members about
their family court proceedings, even if the family members will never
need to serve as witnesses either for the mother or as “hostile” witnesses,
called by the other side. Enjoining a desperate mother from seeking
moral support from her own family members causes her to feel isolated,
demoralized, and browbeaten. Some mothers become suicidal, while
others have succeeded in taking their own lives. \pfl{}

\prHead{Ending Institutional Violence against Mothers and Children}

Using the power of the family courts to punitively strip mothers of cus-
tody of their children, to force adverse premature weaning of infants
from their mother’s breastmilk, and to jail mothers for no reason at all
signifies a reprehensible human rights crisis and flagrant violations of
constitutional, statutory, and case law. The senseless infliction of “insti-
tutional violence” upon unsuspecting women and children demands a
thorough investigation by the United States Attorney General. And this,
in fact, may be more doable than we think. \pfl{}


Legal scholar and first champion of the protective mother, Jeremiah
B. McKenna, to whom we have dedicated this book, researched
the apposite federal statutes to permit a Justice Department inves-
tigation into family courts and their collaborating auxiliaries (Child
Protective Services, law guardians, 10 visitation centers, etc.). Neustein’s
thirty-one-page memorandum to United States Attorney General Mer-
rick Garland, referenced in this book, is premised on those federal
statutes researched by McKenna, which give jurisdictional authority to
federal prosecutors to open up an investigation nationwide—addressing
\notePage{xli}
first and foremost the counties where the most troubling cases tend to
cluster. 11 In essence, what a federal prosecution affords is immediate re-
lief. It signals to the courts that its corrupt practices—where evidence
is deliberately suppressed; records are falsified; and mothers, and even
more so women of color, are threatened, blackmailed, coerced, and
worse—must all come to an end lest the judges and their accomplices
(law guardians, court-appointed mental health experts, child protective
service agencies, and visitation supervisors) find themselves someday in
the same wretched position as the mothers that have come before them. \pfl{}


Accordingly, if judges and their auxiliaries are successfully prosecuted
for organizing, fostering, and facilitating a racket (as will be described
in painstaking detail) this will restore balance and fairness to the fam-
ily court system—a state of natural homeostasis reached after a purging
of all its rottenness. In fact, we may never go back to the family court
paradigm for handling child custody in the first place. It may, given
its deep-seated venality, become defunct altogether, to be replaced by a
panel of concerned professionals who are not driven by any agenda other
than protecting the mother–child bond and shielding the child from
predatory harm, sans the biases that imperil the family courts today. \pfl{}


We will then have in place of a corrupt family court a true “King
Solomon Court.” And the rampant institutional violence against women
suffusing the domestic court system today shall exist no longer. \pfl{}



\prPart{PROTECTIVE MOTHER MOVEMENT — THEN AND NOW}


\prHead{Protective Mother Movement Makes Great Strides in Recent Years}

In the nearly two decades since the publication of the first edition of this
book—which showed how courts consistently punish mothers with the
loss of custody of their children when they make good faith allegations
of sex abuse against the child’s father—we have witnessed a burgeoning
of a social movement to reform the family courts, in part due to the elab-
orate networking offered by social media and, more recently, by Zoom
\notePage{xlii}
meetings that bring together protective mothers across the country for
an informal virtual “get together.” There are now scores of trusted web-
sites (and blogs) posting the latest research studies, news articles, and
video and YouTube recordings, all available to mothers who have lost
their children to the family courts. \pfl{}


Hence, no mother is fighting alone today—isolated in her despair—
for she can now readily avail herself of the multitude of websites, where
some even provide referrals to family law attorneys, expert witnesses,
or give advice on legal strategizing to win back custody of children.
Some propose specific legislative reforms to improve the functioning
of the family courts while others propose a jettisoning of the existing
court system altogether, to be replaced by a more humane structure for
determining children’s best interests. Regardless of their approach, the
mantra of all these sites—National Safe Parents Organization, 12 Moth-
ers ReVolution, 13 Custody Peace, 14 The Women’s Coalition, 15 Safe Kids
International, 16 One Mom’s Battle, 17 The Leadership Council, 18 and Cal-
ifornia Protective Parents Association, 19 just to name a few—is to raise
public awareness about the catastrophically high numbers of mothers
losing custody to abusive ex-partners or spouses. \pfl{}


The term “protective mother”—which has become a part of the ver-
nacular of family law practitioners and women’s groups across the
country—was coined around 1986 by mental health professionals who
had formed the American Professional Society on the Abuse of Children
(APSAC). It refers to a mother who mounts a fervent struggle to protect
her child from all forms of abuse, particularly sexual abuse, commit-
ted by the father, in a hostile family court system intent on beating her
down every step of the way. And part of that struggle for the protective
mother, in addition to mounting a strong defense in her own case, is her
engagement in organizing, assembling, and protesting these inequities
as a collective issue. Some of those activities have included participation
in peaceful demonstrations in front of state capitals around the country
as well as at the nation’s capital. Avid letter-writing/e-mail campaigns,
petitions, and social media posting/likes have served as another way of
voicing one’s concerns about the crisis in the family courts. All in all,
a five-minute Google search of the protective parent/family court issue
\notePage{xliii}
will provide an excellent gauge of the rhythm and beat of this growing
social movement. \pfl{}


\prHead{Plight of Protective Mothers Heeded in Academic Settings}

The protective mother issue now resonates in the halls of the universi-
ties, whereas two decades ago when the first edition of this book was
being prepared for publication there had only been sporadic interest at
the universities in this topic. For example, one of the authors was invited
in the fall of 1990 to lecture at the Women’s Law Forum at Rutgers State
University School of Law 20 only because the professor organizing this
effort, H. Joan Pennington, was a remarkable activist and visionary, hav-
ing founded the National Center for Protective Parents, Inc., as a legal
resource for mothers caught in hostile family court litigation. But such
efforts to bring this subject to undergraduate and graduate classrooms
have been rare. \pfl{}


Today, the Battered Mothers Custody Conference, co-chaired by Pro-
fessor Mo Therese Hannah, is held annually at Siena College in Albany,
New York, where Professor Hannah serves as a full professor of psy-
chology. Her conference has grown exponentially since its inception in
2004; it now serves as the premier nexus for the protective mother move-
ment. More recently, the California Protective Parent Association has
partnered with the University of California Irvine (UCI) Initiative to
End Family Violence; and in so doing, this grass roots parents’ group
has joined a major academic institution to collaborate on an annual,
full-day, accredited continuing legal education symposia (although at
present, held virtually), which had started in fall 2021. The conference,
titled “Forward Together: Multidisciplinary Perspectives on Protecting
Children from Abuse,” 21 had a panel of speakers consisting of several ac-
tivists who addressed the colossal mistreatment of mothers in the family
court system. This points to a sea change in attitude. \pfl{}


That is, whereas mothers’ disenfranchisement in the family courts—
which had at one time been solely the province of activists—had now
entered mainstream institutions for discussion and debate. And as
\notePage{xliv}
previously discussed, the discourse on this topic was not a flash in the
pan; instead, it would become a regularly occurring annual event. This is
a marked change from 2005 when the first edition came out. Back then,
one of the authors had been invited to serve as visiting faculty at the
Tristate Family Law Conference held by the National Judicial College
in Bretton Woods, New Hampshire. 22 The attendees at the conference,
and the conference organizer herself, looked upon this subject as novel,
but no less worthwhile to explore. This represented a healthy discourse
on family courts’ betrayal of women and children inside the protected
corridors of a judicial education institution, although back then such a
forum was more or less an isolated event. \pfl{}


As the years moved forward, and as academic institutions began
to show a continuing and steadfast interest in the atrocities of fam-
ily court, George Washington University Law School would become
another very important venue for broadcasting these profound injus-
tices toward mothers who tried to protect their children from abuse.
George Washington University Law School Clinical Law Professor,
Joan S. Meier, became the Founding Director of the National Family
Violence Law Center (NFVLC) situated at George Washington Univer-
sity Law School. The Center’s function is to use “empirical research,
policy development, and collaborative advocacy to educate profession-
als, the media, and the public, and to advance critical legal reforms
and systemic change in courts’ responses to families with a history of
abuse.” 23 \pfl{}


In fall 2021, Family and Intimate Partner Violence Quarterly published
a study, titled “Protective Mothers, Endangered Children,” spearheaded
by California State University at San Bernardino professor Geraldine
Butts Stahly together with two highly respected activist organizations,
the California Protective Parents Association and Child Abuse Solu-
tions, Inc. The data collected for this study came from a hard-copy
questionnaire, consisting of 101 questions, which had been distributed
at national and international conferences concerning domestic violence
and child abuse. The purpose of this study was to learn about the “ex-
perience of mothers attempting to protect their children from abuse.” 24
The authors of this study were so terribly alarmed by the results began
\notePage{xlv}
to circulate its findings (399 mothers from across the country had filled
out the questionnaire) at multiple workshops and poster sessions ahead
of its publication as they felt an urgency to alert professionals to the
serious failures of the family courts. \pfl{}


Specifically, the study authors witnessed repetitive patterns of moth-
ers losing custody and children being placed at significant risk of harm
by the other parent, although the rates of such system failure were a
bit higher in California, where mothers had been found to lose cus-
tody of their children to identified abusers 85.4 percent of the time as
compared to 74.6 percent in other states respectively. The study authors
asked for an overhaul of every facet of the system—custody evaluators,
mediators, judges, child welfare services, attorneys for children, state law
enforcement—including setting up “a citizen review ‘second look’ com-
mission . . . at the state and federal level to review cases in which children
are currently at risk, with the power to intervene [in order] to protect
[children].” The study concluded with an exhortation for “[i]mmediate
and effective changes [that] must be made to law, policy and procedure
to provide safety for children of divorce and separation.” 25 \pfl{}


\prHead{Protective Mothers Finally Find a Sympathetic Ear in the Mass Media}

On the media front, likewise, we have witnessed significant progress in
addressing the protective mother issue since we first wrote about it in
2005. Those changes in the news media’s posture deserve mention. To
wit, in the early days of the protective mother movement, the media
coverage was sporadic at best. But even more troubling was the fact that
it was so unpredictable at times, often prejudicial to the mother. \pfl{}


As a result, many mothers were “duped” when reporters assured them
the hard news coverage or feature story would be in their favor, only
to learn (sadly, the hard way) that there was an ostensibly robust fa-
thers’ rights agenda massaging the news industry at that time. It was,
therefore, not uncommon to see protective mothers totally flummoxed
when, after telling their heart-wrenching stories to the press (sometimes
working with the reporter for many months as these stories can be very
\notePage{xlvi}
complicated), they would find that the news coverage (both print and
broadcast) of their tragic saga had portrayed them in a terribly unflat-
tering light. In contrast, their ex-husbands or ex-partners, and certainly
the father’s advocates, such as the fathers’ rights organizations, were all
made to look respectable and credible when the facts of the case would
significantly prove otherwise. \pfl{}


More often than not, the copious documentation supplied by the
mothers to prove their innocence had no effect on the outcome of the
news media story. It was therefore not uncommon to see mothers drawn
and quartered by a media world that luxuriated in suffocating the truth.
Thankfully, that horrid reality has changed. Hence, in the years since the
publication of the first edition, we have seen media deliver more reliable
news reports, rather than habitually revert to myths, factoids, and far-
cical prejudices aimed at excoriating protective mothers for no reason
at all. \pfl{}


For example, in 2016, investigative journalist Laurie Udesky, work-
ing for 100Reporters, a nonprofit news organization giving a platform
for freelance journalists writing about corruption and social injustices,
had a feature story appear on the tragic travesties of justice in the fam-
ily courts. Her article was titled “Custody in crisis: How family courts
nationwide put children in danger.” 26 A year later she would be given
the Excellence in Journalism Awards, from the Society of Professional
Journalists, Northern California Chapter, for her work in exposing the
ill-fated decisions of family courts that place children in harm’s way. \pfl{}


In 2019, The Washington Post would write a feature story on the ab-
ject plight of mothers punished with the loss of custody when they
try to protect their children from abuse perpetrated by the other par-
ent. The article, aptly titled “A gendered trap,” would quote Meier,
who found a pattern of judges repeatedly placing children with the
abusive parent: “Meier called the pattern a ‘complete betrayal of the
mission’ of the courts, which are supposed to prioritize the welfare of
children.” 27 \pfl{}


What prompted that prominent news coverage in The Washington
Post was Meier’s large-scale study that found that “mothers who report
abuse—particularly child abuse—are losing child custody at staggering
\notePage{xlvii}
rates.” In fact, the study showed that in over 73 percent of the time when
mothers alleged abuse, and the other parent cross-claimed “parental
alienation” (the notion that when a child refuses to have a relationship
with a parent it is due to purported “manipulation” by the other par-
ent), the family court judges will switch custody to the alleged abuser. 28
The Meier study, as a matter of fact, was able to hold up admirably to
specious attempts made by two critics—with a clear agenda—seeking to
negate these significant study findings. \pfl{}


Namely, in 2022, the Journal of Family Trauma, Child Custody, and
Child Development (formerly the Journal of Child Custody) published a
strong and persuasive commentary by Meier and colleagues that pointed
out that opponents of her study (although only two and certainly not
a huge number) had based their criticism on “incorrect speculations
and misstatements about [their] methodology” and that any possible
“methodological flaws” that may have “legitimately cast doubt on the
validity of the FCO [Family Court Outcome] study’s findings have not,
to date, been identified.” In responding to her critics, Meier and her co-
authors put this matter finally to rest by showing the “lack of substance”
to their critics’ findings “while substantiating the credibility” of their as-
siduous research of family court custody (case) outcome when parental
alienation is raised to discredit the mother. 29 \pfl{}


In her successful rebuttal of her critics’ fallacious arguments, Meier
preserved the important contributions of her study and, no less, the
eye-opening lead story in The Washington Post. This was crucial, as any
progress made in getting exposure for the protective mother issue in a
major outlet, such as The Washington Post, might have been reversed
if the critics of her study went unanswered. What we see here is that
the deeply rooted gender biases casting aspersions on the credibility of
mothers cannot be left unchallenged in either the professional circles 30
or in the mass media. As for the latter, because it is axiomatic that public
opinion drives policy makers and legislators, it is imperative to have me-
dia coverage supportive of protective mothers instead of working against
them. \pfl{}


While the 2019 Washington Post article still remains the most highly
cited news story, today about the undeserved backlash against mothers
\notePage{xlviii}
who in good faith report child sexual abuse to the courts, a handful of
mainstream media venues have, in subsequent years, reported on this
topic as well. For example, in May 2021, NBC aired a two-part series
on children who were murdered by their fathers during unsupervised
contact with their children. Pointing to the lack of judicial oversight or
accountability, child advocates who were interviewed for this NBC se-
ries urged for an overhaul of the court system, including the passage
of laws that would compel courts to consider abuse evidence—both
present and past—when making custody/visitation decisions so as to
ensure children are protected from violence and abuse and, in the most
extreme case, death. 31 \pfl{}


Two years later, in 2023, New Hampshire investigative journalist
and author Olivia Gentile wrote an outstanding feature-length story
for Insider on how the discredited “Parental Alienation” theories have
contaminated the American family courts, causing abused children to
be confined to the custody of their abuser with little or no access to
their mother, who strenuously tried to protect them. Giving as example
a San Diego case, Gentile showed how a mother lost custody of her
children after her son’s therapist reported to the court how the twelve-
year-old boy was forcibly coerced by the father to perform lewd acts
that the child described as “disgusting.” In her reporting on this story,
Gentile showed that several years after the mother lost custody of her
three children to their abusive father, the mother would serendipi-
tously find in a cloud storage account she had once shared with her
ex-husband, numerous sexually explicit images (both video and still
shots) the father had taken of the children after he obtained custody of
them. 32 \pfl{}


The effects of the sympathetic journalists resonated with the mothers.
In our interviews with women who had conversations with reporters,
we found that journalists like Gentile were deeply appreciated by the
protective mothers. What we heard was that she gave generously of
her time to distressed mothers, often answering their e-mails—albeit
at late hours of the night, and on weekends and holidays, too. Notably
these reporters were not feminist writers with a long history of champi-
oning women’s causes; instead, they were sensitive and caring journalists
\notePage{xlix}
who witnessed an injustice meted out to protective mothers and ap-
plied their journalistic skills to bring these family court travesties to
light. \pfl{}


\prHead{Feminist Writers and Researchers Come Together to Champion
the Plight of Protective Mothers to a Broader Audience, Outside
of the Traditional Feminist News Outlets}

This past decade would prove to be historic for broadcasting the pro-
tective mother issue to a wider readership. Whereas in the early days,
namely the 1980s, 1990s, and early 2000s, it was mainly feminist newslet-
ters and magazines that had profiled the wretched fate of mothers
peregrinating through the family court and child welfare system, today
one can find discussion of these critical issues in venues outside the fem-
inist media. And the researchers who have labored to study this problem
have much to do with the progress in that regard. \pfl{}


For example, Providence Journal columnist Anne Grant, who is the
former executive director of the Women’s Center of Rhode Island, has
paid much attention to the gender biases pervasive in the family courts
in her writings on social inequities and unbalanced power across vari-
ous sectors of society. The Providence Journal is certainly not a feminist
publication as it is a general-purpose newspaper widely read in the New
England area. In her Amazon review of former senior social science an-
alyst at the National Institute of Justice Leora Rosen’s book, Beyond the
Hostage Child: Toward Empowering Protective Parents, 33 Anne would
write:

\prQuote{This book holds validation for those who have been traumatized
when courts removed terrified children from protective parents
and gave them to the sole custody of abusers. Dr. Rosen shines
a light we need to go forward. \quasiPar{}
She asserts that alleged crimes of domestic violence and child
sexual abuse within the family should never be sent to civil
courts that are designed for compromise. She briefly describes five
\notePage{l}
proposed models for change and offers more detail on a sixth,
composite model, Child At Risk Classification Office (CARCO)
that focuses on a public health assessment of the child's risk of be-
ing exposed to violence or abuse. She uses the acronym TRIAL to
represent key elements of CARCO: Training, Reporting, Investi-
gation, Adjudication, and Long-term planning—that are woefully
absent from the present practice of adversarial litigation in family
court.}

\noindent{}Grant concludes with plaudits for the author’s vision, dedication, and
commonsense approach to the protective mother debacle:

\prQuote{Dr. Rosen has performed a huge service by focusing on those of us
who feel numbed by our own inability to protect desperate children and
non-offending parents from the lies of lawyers and psychologists who
have reduced them to a profit center. She [the author] concludes by urg-
ing Congress to use its authority and enact CARCO for the District of
Columbia, creating a model for the nation. Federal funding incentives
can be redirected to inspire other states to follow suit and to end the
nightmare that breeds child abuse at family court.} 34 \pfl{}


\prHead{Protective Mother Issue Gains Stardom in TV Miniseries
now Streamed on Hulu}

The media coverage of the protective mother’s plight would mushroom
in recent years and prove worthy of a Disney Channel platform. By
summer 2022, FX would air a five-part docuseries, “Children of the
Underground,” which would be subsequently streamed onto Hulu,
where it remains to date. The Emmy Award-winning production com-
pany, Story Syndicate, had diligently put over three years into the
research and production of this spellbinding television miniseries. Ex-
ecutive Producer Ted Gesing would contact us in the fall of 2019, after
reading the first edition of this book, nonplused and horrified that the
family courts had descended into such a state of damnation. The docu-
mentary convincingly showed the angst of mothers and the cruelties of
\notePage{li}
the family court system that oppress mothers to the point where many
had no choice but to take their children and go on the run—and suffer
poverty, lack of medical care, isolation, food deprivation, and chronic
stress as a result. \pfl{}


This docuseries on the protective mother issue received good reviews
in a number of different media venues: New York Post, The Daily Beast,
A\&E’s True Crime Blog, and Indiewire, just to name a few. It was praised
on www.rogerebert.com for “captur[ing] the growing pains of under-
standing mental health, which includes the people and children who
were mishandled in the [court] process . . . [and is] all the more sober-
ing when one realizes how certain parts of the family court have not
changed, or that the same problems can arise.” 35 \pfl{}


Yet, some protective mothers and advocates had mixed feelings
about the docuseries. They did not take issue with its content. On the
contrary, they found this docuseries—in which mothers and (now)
grown children had so poignantly and persuasively bared their souls—
very compelling. However, they had thought the gravity of the injus-
tice and eloquence in the way the story had been told would finally
catapult this issue to the front burner. That is, they perceived the
series as mainly catalytic in stirring public outrage over the shame-
ful abuses of the family courts and they wondered why that hadn’t
happened. \pfl{}


However, looking at this situation more objectively, who can say that
the docuseries won’t have that effect; perhaps if not today, then tomor-
row? As with any social movement, the exact point of critical mass is
hard to predict and, similarly, just as hard to identify when it actu-
ally does occur. In fact, it seems that historians are much better able to
time-stamp critical developments in any political or social movement
than those living through that epoch. For this reason, it would be fool-
hardy to throw in the towel just because the walls of Jericho didn’t come
tumbling down in the wake of the docuseries. In fact, much traction has
been built up already. Who would have imagined back in 2005 when this
book was first published that someday the topic would occupy a five-
part television miniseries airing on a major network and subsequently
streamed onto a Walt Disney platform? 36 \pfl{}

\notePage{lii}

\prHead{Women’s eNews Recognizes the Protective Mother Issue}

Under the direction of Executive Director and Editor-in-Chief, Lori
Sokol, PhD, Women’s eNews, the award-winning nonprofit news ser-
vice that covers issues of concern to women in the United States and
around the world, has welcomed editorials on the protective mother is-
sue. Publishing on this news service has been a real boon to protective
mothers because Women’s eNews is a well-respected global forum for
women to share their views on public policy and social reform. \pfl{}


In February 2023, best-selling novelist Talia Carner wrote a com-
pelling piece for Women’s eNews on why protective mothers, in some
cases, may be perfectly justified in defying family court orders. 37 Around
the same time, Women’s eNews published Amy Neustein’s editorial on
the looming civil rights crisis in the family courts. Neustein described
how mothers are routinely deprived of due process and equal protection
under the laws and, perhaps most egregious, they are denied their civil
rights to self-determination to feed their infants from their own breast.
She pointed to a common pattern of mothers being court-ordered, un-
der threat of jail and fines, to expel their breastmilk into a bottle so
that the father can have generous and unrestricted long weekends and
overnights alone with the infant. 38 In many of these cases, the fathers
had been absent during the pregnancy and birth. They do not pay child
support either. Yet, they demand the infant spend solitary time with
them, even though it disrupts the infant’s nursing schedule and the
mother–child bond. \pfl{}


Women’s eNews would welcome a number of subsequent editorials
from Neustein. The pieces looked at family court from various an-
gles. These included a post-mortem examination of a high profile case
wherein a New York mother, herself an attorney and former federal
prosecutor, took her own life after years of traumatic abuse by the
Westchester Family Court; the harsh removal of custody from mothers
who had failed to speak eloquently on the witness stand or had become
emotional in court upon facing an abusive and controlling ex-spouse
or partner who was trying to take their children away from them; the
spontaneous formation of the “Mothers’ Vote,” an ad hoc voting bloc


\notePage{liii}
of protective mothers across blue states, red states, and swing states
who are united in their family court plight; and the necessity of go-
ing outside the judicial branch to another branch of government to seek
remedy when the courts as an entirety have failed to protect women and
children. 39 \pfl{}


We noted in our first edition that we had spoken at length to Rita
Henley Jensen, then executive director of Women’s eNews, about the
protective mother issue. Ms. Jensen was extremely sympathetic, point-
ing to the daily calls she was receiving from desperate mothers losing
their children in family court. However, since that was twenty years
ago, the protective mother movement had not yet acquired a unified
voice. Few, if any, articles had been published in Women’s eNews on
the protective mother plight. Today, the situation is much different. This
preeminent digital news service has emerged as the premier platform for
sharing in depth commentaries on the family court system. Certainly,
their wide readership in government and in the private sector will gain
insight into what mothers have been experiencing on a daily basis in the
American family court system. \pfl{}


\prHead{Protective Mother Movement Builds up Traction
on the Political Front}


As an example of the kind of traction built up on the political front,
in 2017, a group of activists from across the country, led by Constance
Valentine, founder of the California Protective Parents Association,
brought one hundred copies of the first edition of From Madness to
Mutiny down to Washington, DC, 40 and left copies on the desks of the
crucial members of the House of Representatives. Soon thereafter, Ms.
Valentine and her co-activists were successful in getting the House to
pass a resolution to prioritize child safety in custody and visitation cases.
“H.Con.Res. 72,” as it would officially be called, was sponsored by Mary-
land Congressman Jamie Raskin, along with several of his colleagues,
who recognized the urgency of putting forth a resolution to affirm the
importance of child safety during family court litigation. 41 \pfl{}

\notePage{liv}
Although the Resolution passed in 2017–2018, it left open the possi-
bility of taking the matter, albeit sometime later, to its next level. Written
into the Resolution was a suggestion/provision for a full Congressional
investigation of the family courts that have been proven to show a men-
acing, baleful and sinister effect on the welfare of children because they
habitually turn innocent children over to the abusive parent, while forc-
ing the non-offending mother to stand by and watch her children suffer,
and sometimes die, from heinous abuse or neglect. Indeed, when that
next step is reached, the family court backlash against protective moth-
ers will finally be given a front-row seat on Capitol Hill—and lawmakers
will listen attentively. \pfl{}


In March 2022, President Biden signed into law the “Keeping Chil-
dren Safe From Family Violence Act” (commonly referred to as “Kay-
den’s Law,” named after a seven-year-old Bucks County girl (Kayden
Mancuso) murdered by her abusive father during a court-ordered un-
supervised visit with the child. Meier, Pollack, and colleagues produced
assiduous research to impress upon Congress the need for a special pro-
vision to be included in the reauthorization of VAWA (Violence Against
Women Act) so as to protect vulnerable children from injury and death
at the hands of the abusive father. Seven states, led by Colorado as the
flagship state, 42 have already enacted Kayden’s Law, in whole or in part.
They include Colorado, Utah, California, Tennessee, Maryland, Penn-
sylvania, 43 and Arizona. Connecticut is now close to enactment as well.
Whilst many more states need to join this effort, the states that have
taken the lead clearly demonstrate the power and force of federal legisla-
tion in inspiring states to enact laws to protect children in custody battles
from being subjected to injury and death at the hands of a parent with
a known history of spousal abuse and domestic violence, particularly
when tied to federal funding. \pfl{}


While the effects of VAWA have evidently trickled down to the state
level, individual states have, similarly, on their own initiated legisla-
tion in response to tragic deaths of children in their state. For example,
in 2023, several New York senators co-sponsored “Kyra’s Law,” which
was named after a two-year-old Long Island girl forced under court or-
der to have an extended visit with her violent father in Virginia who
\notePage{lv}
fatally shot her while she slept. “Kyra’s Law” emphasizes the importance
of giving sufficient evidentiary weight to claims of domestic violence
when determining the custodial and visitation arrangements for a child.
The bill stresses that keeping a child safe from a parent with a history
of violence should be the first priority of the family courts. Although
many of these state bills, such as Kyra’s Law, have not yet been able to
gain the full support of both state-legislative houses needed for pas-
sage into law, the passionate commitment of a handful of outstanding
legislators has, nevertheless, been propitious for committee passage,
which gives reason for hope that such bills will eventually make it to the
floor. \pfl{}


On the whole, this second edition of From Madness to Mutiny is being
written at an auspicious juncture in women’s history. We are witnessing
a collective uprising as mothers, advocates, and scholars unanimously
urge for an end to violence (both personal and institutional) against
women and children in the family courts across the globe. As demon-
strated in the previously discussed United Nations Human Rights Office
of the High Commissioner study, the family courts have become em-
blematic of the ghastly violence against women and children at every
phase of the judicial process. And the mothers are determined to no
longer stand silent in the face of such atrocities—both now and in the
future. \pfl{}


\prPart{SEVERE MISTREATMENT OF BLACK, INDIGENOUS, AND PEOPLE OF COLOR MOTHERS IN THE FAMILY COURT}

\prHead{PROTECTIVE MOTHERS SUFFER TWO FOLD IN THE FAMILY COURTS}


In the new material provided as supplements to a number of the chapters
from the first edition, we examine in detail the mistreatment meted
out to Black, Indigenous, and People of Color (BIPOC) mothers in the
family courts. Although we had presented some compelling material in
the first edition showing, for example, the jailing of a BIPOC mother
from Brooklyn at the caprice of the family court judge, we hadn’t ex-
amined the family courts and the child welfare system from a racial
\notePage{lvi}
perspective per se. Now that we have, we can show that in the closed and
claustrophobic world of family courts racism flourishes and is given new
life, adding undoubtedly to the execrable mistreatment already meted
out to protective mothers. \pfl{}


\prHead{Black, Indigenous, and People of Color Mothers Suffer Racial
Discrimination on Several Fronts}

First, BIPOC mothers are slandered by name-calling in the family
courts. In addition to the routine psychiatric name-calling directed
at protective mothers (“disturbed,” suffering from “personality disor-
ders,” “narcissistic,” “grandiose,” etc.) women of color shoulder the extra
burden of having racial epithets hurled at them. For example, in the
Brooklyn Family Court, a Latino mother, who was also a police of-
ficer, pleaded with the family court judge for the protection of her
five-year-old girl from acts of sodomy at the hands of the father. The
father had forced himself on his daughter during an overnight visita-
tion and the child came back hysterical, crying how she “almost choked.”
But when the mother made entreaties to Brooklyn Family Court Judge
Leon S. Deutsch (now deceased) to protect her child, he became en-
raged. Bellowing at the mother, he called her “an hysterical Hispanic!”
Although the term “Hispanic” is not much used anymore, Deutsch was
an older male, Caucasian judge, whose vocabulary reflected an earlier
era and his cultural background. As such, he drew from his arsenal of
derogatory and discriminatory language to assail this protective mother
who had no other agenda in mind than to protect her young, helpless
daughter from the brutalities of sexual assault and sodomy at the hands
of her abusive father. \pfl{}


Second, BIPOC mothers are more likely to be incarcerated for con-
tempt of court than their Caucasian counterparts, albeit on unsubstanti-
ated grounds. Interestingly, when comparing the incidence rate of mere
“threats” of contempt of court (as opposed to acting on those threats
by jailing mothers), we did not find a higher incidence of threats made
against the BIPOC mother. For example, both demographic groups of
\notePage{lvii}
mothers are equally threatened by family court judges with contempt of
court if they would so much as “speak” to anyone (and sometimes this
includes parents and/or siblings) about their family court proceeding. 44
But, when it comes to actually jailing mothers, BIPOC mothers are in-
carcerated more often than Caucasian mothers, which is very much in
keeping with the stats across the general population that show a much
higher percentage of incarceration among people of color. \pfl{}


Third, when family courts strip a BIPOC mother of custody and then
orders supervised visitation between the mother and her child (“super-
vision” is used so as to prevent the mother and child from speaking
about abuse, sexual and otherwise, committed by the father, and it also
functions as a hedge against the child who may initiate discussion in
making new disclosures of abuse), the mother is often forced to see her
child at the local police station, instead of at a court-contracting visi-
tation setting, as described in the next section. Interestingly enough, a
judge’s ordering of supervised visits at the police station is not always
contingent on whether or not the BIPOC mother would have the fi-
nances to pay the visitation supervision fees held in a private setting. In
studying the family court backlash against protective mothers for nearly
four decades, we have never ever seen Caucasian mothers ordered into
a police setting to see their children. Such arrangements are appalling,
and they naturally have injurious effects on BIPOC mothers and their
children. \pfl{}


When children see their mothers in the police station, they naturally
associate trepidation and fear with those visits. As they sit in this sterile
setting, surrounded by men and women in uniform—clad with gun-
filled holsters and handcuffs—they often cannot relax on the visit with
their mother. BIPOC families have grown to fear police because they
are often the target of police harassment (sometimes tragically leading
to injury and death), whereas Caucasian families are more likely to view
police as people to whom they can appeal for help. Forcing a BIPOC
mother to see her child at a police station revives the long-standing
trauma for people of color who try their very best to avoid contact with
police altogether, and for good reason. Moreover, a child seeing that
their mother must be supervised in a police setting will understandably
think their mother did something terribly wrong, otherwise why would
\notePage{lviii}
she need to be monitored at a police station? In short, in as much as
it is common for trepidations about arrest to constantly hover over the
BIPOC population, arranging visits for the mother at a police station is
wholly destructive to the mother–child relationship. Any angst that the
mother feels sitting in a police station can be transferred to the child,
making the visits uncomfortable (and often traumatic) for both mother
and child. \pfl{}


We followed a case in which a Bronx Family Court judge had ordered
the mother to have her visits with her four-year-old daughter (who was
sexually abused by the father) at the police station. The mother had
presented medical evidence to the court to substantiate the abuse, but
the judge refused to listen. The child’s pediatrician had backed up the
credible evidence of abuse. And although the child herself gave narra-
tive accounts of how she was defiled, neither the law guardian nor the
judge heeded the mother’s pleas to protect her child. The judge, instead,
stripped the mother of custody and transferred the custody of the very
young girl to the father. The mother held a secretarial position at a For-
tune 500 company. So, she could have paid for visitation supervision at a
private setting, but the judge forced this BIPOC mother to see her child
at the police station instead. This was the only venue she would allow
for the mother to have visitation with her child. The mother remarried,
a lawyer in fact, and continued to fight for her child. \pfl{}


Sadly, her new husband, despite all his motions to the Appellate Divi-
sion, was unsuccessful in vacating the draconian visitation arrangement
at the police station, as well as in overturning the custody award to the
sexually abusive father. The mother, worn down and defeated by these
horrid conditions, was unable to continue her life as a middle-class wife
married to a young, aspiring lawyer. Her marriage fell apart and she
moved in with a relative down South. Her daughter remained up North.
When we asked her husband about whether racism had played a role in
the case, he didn’t hesitate for a moment. He exclaimed, “For sure . . .
and throughout the entire case!” \pfl{}


Fourth, whereas family courts are weaponized against protective
mothers, those who are BIPOC are sometimes precluded from the court
process altogether, and that can leave the mother with no redress at all.
For example, a Black special education teacher living in Georgia was
\notePage{lix}
told last year that she could not even have a hearing before a judge be-
cause there are “no available transcribers.” This was untrue—the same
Georgia court had made itself available to hear cases brought by Cau-
casian mothers. Instead, this teacher was shunted to a mediator in lieu of
a judge, as the court fell back on its lame excuse that there were no court
reporters available. The mediator hastily rushed the case of this Geor-
gia mother, sending her sexually abused seven-year-old daughter right
into the hands of an abuser, and subsequently terminated the mother’s
visitation privileges with her daughter. \pfl{}


The mediator made this rash decision to cut off the mother’s visits
with her daughter purely based on hearsay evidence provided by the fa-
ther. The mother could not rebut this evidence because she had been
denied access to the courts. A year later, the mother has still been de-
nied a hearing before a judge, and she still cannot see her child. The
child suffers from a congenital disease that has required multiple hospi-
talizations. The mother, who had been at her daughter’s side throughout
all these hospitalizations, was now barred from seeing her daughter
because the mediator cut off her visits based on the hearsay evidence
supplied by the father. Most tragically, when the visits were cut off, the
child would soon need another hospitalization. Admitted to Philadel-
phia Children’s Hospital to receive blood transfusions and dialysis, the
BIPOC mother was forbidden to come to the hospital to see her daugh-
ter. Speaking to one of the authors, she justifiably said she feared “she’d
never see her daughter again on earth.” \pfl{}

\prPart{UNSAVORY “BUSINESS MODEL” TAKES HOLD IN THE FAMILY COURTS}

\prHead{The New Economic Realities of the Family Courts—and How
It Affects Protective Mothers}

The new business paradigm found in the family court system today
forces mothers into accepting services at court-contracting agencies,
such as those that provide visitation supervision. For context, moth-
ers have traditionally been asked to pay for services ordered by the
\notePage{lx}
court. For example, judges for years have been ordering mothers to pay
for a court-appointed custody evaluation, costing on average \$15,000
for the evaluation and the ensuing report made to the court. Today,
however, mothers must pay, in addition to the custody evaluation, fees
for the visitation center facility to monitor their visits with their chil-
dren. Often paired together with the costs of supervised visitation is
the judge’s mandate that, as a precondition for visitation privileges, the
mother must go into “therapy” with a court-approved mental health
practitioner, whose sole purpose is to disabuse the mother of their “mis-
perceptions/delusions” about the abuse their children have suffered,
even in the face of credible evidence, at the hands of the father. \pfl{}


Often these costs are so prohibitively expensive to the mother that
she will stipulate to custody and visitation arrangements—undoubtedly
prejudicial to her—merely to be relieved of paying such costly visita-
tion supervision fees that she can ill afford. Such a set up potentiates
extortion of the mother’s rights to custody of her child as she will be
forced to “agree” to unfair custody terms that will place her child at
risk just to relieve herself of the onerous financial burden of paying for
weekly visitation supervision. As a matter of fact, the way that many
family court visitation programs legitimate their charging of protective
mothers is that they claim to maintain a (discretionary) sliding scale.
However, in practice, while these visitations centers subsidize the costs
for low-income families, such as in cases of inadequate guardianship
where children go hungry, miss school, and receive no medical care,
they charge the middle-class protective mother at full rate. \pfl{}

\prHead{Court-Ordered Visitation Used to Prevent the Child’s Disclosures
from Investigation}

When the courts order supervised visitation for the mother, the “su-
pervision” is used for two purposes: first, to prevent the mother and
child from speaking about sex abuse committed by the father; and, sec-
ond, to act as a deterrent against the child initiating discussion with the
mother in making new disclosures of abuse. Consequently, as we have
seen over and over again, if the child does make a new disclosure, and
\notePage{lxi}
many of them will because child molestation is a recidivistic crime, the
court-appointed visitation monitor is poised to discredit the mother by
denying that such disclosures were made in the first place. Regrettably,
supervision monitors cannot act independently in writing up such new
disclosures made by the child because the centers receive all their re-
ferrals from the family courts entirely. Consequently, they are chary to
antagonize a judge who has already labeled the mother a “parental alien-
ator” and is hence closed-minded about the abuse. As a whole, such a
court “business model” is, for obvious reasons, deeply flawed and coun-
terproductive at best, and inherently deceitful and venal at worst. Yet,
it flourishes in the family court system. What follows is an elucidation
of just how well these visitation centers are faring financially in today’s
family courts. 45 \pfl{}


\prHead{Visitation Centers Impose Undue Financial Hardship on Mothers}

First, the visitation supervisor often works for just a fraction above
minimum wage (\$16 an hour) to maybe a maximum of \$50 an hour.
Nonetheless, the mother is charged about \$140 to \$160 per hour for the
visit with her child, appropriating the bulk of the hourly fee to the visi-
tation center. Supervised parents are prohibited from engaging relatives
or friends to supervise their visits, which would certainly relieve most
mothers an onerous financial burden. The reason given for this is that
the courts must “monitor” and “regulate” these visits, wherein visitation
supervisors are obligated to write regular progress reports to the court.
If one considers, on average, that mothers will be given four hours of
supervised visitation per week, sometimes divided into two visitation
sessions, the weekly costs add up very quickly. \pfl{}


Second, most court-contracting visitation centers require “orienta-
tion” fees to be paid by the supervised parent up front before the
visits can be arranged by the visitation facility. Orientation fees can
range anywhere from \$150 to \$450, depending on the region of the
country where the visitation centers can be found. During these ori-
entation sessions, the mother is told that her behavior will be closely
\notePage{lxii}
watched. They are cautioned that notes will be taken on their visits
and that the supervisor will be reporting regularly to the judge. What
is at stake is certainly not an appraisal of how well the mother relates
to the child (as no one has ever accused protective mothers of being
deficient in nurturing and loving their children), but whether indeed
she “talks against” the father to the child or whether she gets “an-
gry” at the father if the child should make a new disclosure of sexual
abuse. \pfl{}


\prHead{Visitation Centers Impose Punitive Conditions on Mothers
and Children}

Mothers are not only restricted in what they can say to their children,
but they are also enjoined from normal mother–child activities. For
example, they are warned not to bring any food or drinks during their
visits with their children, as that can create a burden for the visitation
facility in cleaning up afterwards. \pfl{}


In a recent Manhattan Family Court case, the protective mother was
cautioned during her orientation session with the visitation center not
to bring any gifts for her five-year-old girl because the law guardian
and the father’s attorney felt that would be a way of “winning” the
child’s loyalty to the mother, which would have been prejudicial to the
father. In fact, in this case, once the visitation supervisor showed sym-
pathy for the mother and urged the court to allow the child to spend
overnights with the mother, unsupervised of course, the visitation su-
pervisor was removed from the case and replaced by another supervisor,
who was tepid toward the mother at best. Similarly, in a Brooklyn case,
the mother was punished with a suspension of her visits because she bent
her head down and whispered “I love you” in the ear of her precious
six-year-old daughter sitting on her lap. When her visits were re-
stored months later, the judge issued an order forbidding the protective
mother from “whispering” to her child; to enforce that order, the visita-
tion supervisors never allowed the young child to sit on the mother’s lap
again.  \pfl{}

\notePage{lxiii}
\prHead{Mothers Must Pay for Court-ordered “Therapy” to Visit
Their Children}

Unfortunately, the costs to the protective mother do not end here. As
previously mentioned, in addition to paying for her supervised visita-
tion, and for the orientation at the visitation center, the mother must
pay for “therapy” as well. Such court-ordered therapy sessions serve as
a condition that the mother must meet to have visitation with her child.
These are usually weekly sessions and sometimes they are twice a week.
The therapist is then ordered to make regular reports to the court and
to the visitation center on the progress made in those sessions with the
mother—effectively using those sessions to disabuse the mother of the
fact that her child had been sexually abused by the father: in other words,
“gaslighting” the mother until she starts to doubt her own sanity. The av-
erage cost of those “therapy” sessions paid by the mother is \$150 to \$200.
And those therapy sessions must continue for as long as the mother has
visitation with her child. When adding the therapy costs to the visitation
supervision costs, the protective mother is burdened with an average of
\$750 to \$800 per week to see her child. \pfl{}


\prHead{Family Court Arrogated by “Ransom” Business Model Spanning
in Many Directions}


As shown, taking a child away from a mother without justification and
then forcing her to pay to see her child is de facto tantamount to ransom.
And previously discussed, this is a relatively new phenomenon, which
was not as systematic and routinized when we compiled our research
for the first volume. We have detected another scenario in family court
in recent times that likewise follows the ransom paradigm. To wit, chil-
dren are wrenched from their mothers without a hearing and without
any evidence before the court. The mothers are then ordered to pay for
an expensive custody evaluation. Having already lost custody, the moth-
ers are certainly in no position to refuse—what other choice do they
have? \pfl{}

\notePage{lxiv}
This practice of wrenching first and then later performing an
evaluation—in essence, a “fishing expedition” to justify an unwarranted
removal—bears the indicia of ransom. Let’s look at this process more
closely for a moment. \pfl{}


In the past, children primarily stayed with the mother—the primary
caregiver—until the court-appointed custody evaluator submitted their
report to the court. Only after a full trial wherein the evaluator testified
and was cross-examined did the judge then make a custody determi-
nation in the form of a written decision. But today, the courts instead
remove the child from the mother first, and then order the mother to
undergo a court-ordered forensic evaluation. The fees charged for this
court-ordered exam are extravagant because the so-called neutral ex-
perts nesting in the courts depend on these fees for income. Certainly,
had the mother not already lost custody, she may be in a position to
object to such a costly custody evaluation of her parenting skills. But
by taking her child away first, the mother, aching for the return of her
child, will pay any amount to get her child back and that includes a
court-ordered evaluation at a profligate rate. \pfl{}


The court’s appointment of a custody expert only after the mother
had already lost custody begs the question of how the judge can remove
the child from the mother in the first place as no court-ordered eval-
uation report had yet been before the court? Certainly, these are not
mothers with who are active drug users (most don’t even have a history
of drug use), or those with DUIs or arrest records of any kind, which
might be justification to remove a child from its mother even before a
full-fledged parental fitness evaluation was performed. Instead, where is
the evidence before the court to warrant taking away the child from the
mother if the custody evaluation had not yet been performed? \pfl{}


\prHead{Broad Socioeconomic Ramifications of the Ransom Business
Model in the Family Courts}


This “ransom” business model that exists in the court system today, in
which coercion is used to exact payments from mothers after seizing
\notePage{lxv}
their children without statutorily mandated hearings to justify their ac-
tions, has broad economic and social consequences that are important
to note. First, middle-class mothers can ill afford the court ransom fees,
which are heaped on top of the expenses that mothers pay when they
lose custody to the child’s father. They pay, on average, \$2,000 a month
for child support, and to that there are the added costs of healthcare in-
surance and tuition for the child. Besides, the costs to protective mothers
may be even greater at times. \pfl{}


For example, in a Long Island, New York, case, the father accused of
sex abuse won custody of his two young children and then asked that
the mother pay for “childcare services” on top of all the other expenses
she was forced to pay. How ironic: when the children were with the
mother, her parents babysat while she went to work, and they certainly
didn’t charge for their services. But once the mother raised a sex abuse
allegation and lost the custody of her children, her own parents were
then forbidden to see their grandchildren. As a result, they were now
prevented from providing the “free” childcare for the young children
that they had been doing all along. Interestingly, although the children
had a set of paternal grandparents living nearby, the father wanted the
mother to pay out-of-pocket for childcare when he was unable to care for
them. \pfl{}


Childcare expenses, sometimes seen by the courts as a legitimate ex-
pense to level against the protective mother, drain the mother even
further. This example points to the fact that by eviscerating mothers
from the lives of their children, the entire family ecosystem is disrupted:
maternal grandparents who had hitherto provided free childcare so that
their daughters can go to work are now eclipsed from the picture, and
the father must find hired help to stand in place of the childcare provided
by the maternal grandparents. Moreover, can the love of grandparents
be replaced with hired help? \pfl{}


Adding to the disruption of the family ecosystem and family stabil-
ity is the fact that protective mothers bankrupted by the family courts
must depend on their own parents for financial assistance. In prepa-
ration of this second edition, we interviewed many protective mothers
who have explained how their own parents depleted their retirement
\notePage{lxvi}
funds to help them shoulder the enormity of costs incurred by the
ransom business model ruling the family court system today. And if
these mothers were to file for bankruptcy protections, as some of them
have, they are still not protected by bankruptcy laws. This is because
bankruptcy court takes time, and the mothers can ill afford to let weeks,
months, or years slip by while they miss the chance to see their chil-
dren grow and develop. This happened to a Houston, Texas, mother
who filed for bankruptcy but the family court wouldn’t waive the costs
of her supervised visits, essentially rendering her bankruptcy protec-
tion moot. Sadly, she ran out of funds and lost all contact with her
children. \pfl{}


So, the mother is left with little choice: either she gives into the
extortion attempt, paying the ransom to see her children, or suffers es-
trangement from them. Because bankruptcy proceedings do not absolve
the financial burden placed on a mother by the family courts (and this is
for good reason because if one could discharge child support obligations
in bankruptcy proceedings, few fathers would pay support) at the end
of the day the protective mother is faced with a Hobson’s choice. Either
they impose a financial burden on their own parents to exhaust their
retirement accounts, placing them in a dangerously weakened financial
position, or they abstain from asking for help and accept never seeing
their children again. \pfl{}


As we can see, the sociological consequences of a criminally dys-
functional family court system creates immense suffering and financial
insecurity for large segments of the population, most notably retirees
who can ill afford to deplete their nest eggs. What this shows is that fam-
ily court venality has multigenerational repercussions—for the mother,
the child, and the grandparents. It is, therefore, a debacle that engulfs
multiple generations, weakening the family structure and society for no
other reason than to promote an unsavory business model that preys
upon battered women and their abused children. In this setting, moth-
ers become enslaved to the family court’s agenda of protecting lucrative
enterprises that nest in the courts in much the same way that serfs had
been forced to cater to the malevolent self-interests of the nobleman in
medieval times. \pfl{}

\notePage{lxvii}
\prPart{NARRATIVE CASE HISTORIES AND COURT RECORDS: A LONGITUDINAL VERSUS A 
CROSS-SECTIONAL STUDY}

\prHead{How the Second Edition is Organized, and How It is Distinguished
from the First Edition}

Most notably, before discussing the differences between the first edi-
tion and the second edition, we must point out the following detail.
As with the first edition, we have in the second edition, likewise, pri-
oritized the confidentiality of the families whose cases we analyze and
dissect. Thus, we refrain from citing names of parents and children
when presenting narrative data. In addition, when quoting verbatim
from court documents we have redacted the names of parents and chil-
dren. No docket numbers are given either, even though in some states
these records are publicly accessible, because our purpose is to protect
the privacy of parents and children. As with the first edition, we have
copiously documented each case so that, in the event that either con-
gressional leaders and/or the Department of Justice should decide to
pursue an investigation of the scandalous family court system, we will
be able to direct them to the mothers whose cases are profiled in this
book, as well as to other cases in our database that have not been in-
cluded in this volume. We find that the court records, when analyzed in
painstaking detail, as we have tried to do in this book, provide a
sufficiently compelling case for a federal investigation. The harrowing
picture presented in mothers’ narratives, as relayed to us, are well cor-
roborated by the court record itself. This is why we have consistently
pointed to the record as an instantiation of the madness prevailing in
the family courts. Admittedly, although this process of reviewing exten-
sive court files is habitually time-consuming, it nevertheless strengthens
the terrifying narratives provided by the mothers who are haplessly
caught in the vortex of family court litigation. All in all, over the past
five years, in preparation of this second edition, we have listened at-
tentively to mothers’ narratives as part of our sociological study of
the courts, and we have compared those narratives against the case
record. \pfl{}

\notePage{lxviii}

Our approach in the second edition, however, is distinguished from
that of our first edition in several ways. \pfl{}


First, rather than instantiate, as we did in the first edition, the mani-
fold indicia of family court madness and the resultant parental mutiny
by pointing broadly to as many protective mother cases illustrative of
this phenomena as possible, we concentrate instead on fewer exemplary
cases but analyze them in greater detail. We decided on this approach be-
cause our first edition, using a substantial number of cases for illustrative
purposes, had already laid out the themes pervasive in the family courts
(themes which still persist today); therefore, there was no need to repeat
that process. Although for this new edition we hone in on fewer cases,
albeit in much more depth, those we chose to present in this book from
the many that we had reviewed should not be considered an anomaly.
On the contrary, they are truly representative of the egregious abuses
present in the family courts today. The longitudinal cases presented in
this book are, nevertheless, drawn from various regions of the country
so that they are not exclusive to one particular geographic area. \pfl{}


Overall, our presentation of a few cases in depth, as opposed to
broader analyses of multiple cases, allows us to narrate a case longi-
tudinally and show, for example, how child protective services have
become increasingly more aggressive toward protective mothers as their
case progresses and, in particular, increasingly reproachful and excori-
ating toward women of color. Our lengthy case narrative approach also
allows us to show, for example, how law guardians have come to com-
port themselves in such a reckless manner that they arrogantly ignore
well-documented criminal histories of fathers that have a direct bearing
on the safety and well-being of a child. \pfl{}


Second, although in preparation of both our first and second edi-
tions of this work we immersed ourselves in the study of full court
files—transcripts, orders to show cause, forensic reports, motions and
cross-motions, child protective service petitions/case summaries/safety
plans, decisions, appeals, and so forth—for our second edition we had
the opportunity to examine the pleadings from the federal courts as
well because in recent years many more mothers have brought civil
rights actions in federal courts albeit in a futile attempt to seek remedies
\notePage{lxix}
they were not able to obtain in the state courts. Even though most, if
not all, mothers had their cases dismissed because of collateral estop-
pel issues and related “abstention” doctrines (e.g., Younger v. Harris;
Rooker-Feldman; “domestic relations exception”) that prevent federal
courts from intervening when a matter is before a state court, their
federal lawsuit papers nonetheless offered “goldmines” of information,
demonstrating the procedural defects/violations pervading their family
court cases. \pfl{}


Third, in researching protective mother cases, both for this new edi-
tion and for the prior one, we employed a qualitative (rather than a
quantitative) research method, our goal for this current work was to
use close analytical inspection to identify “patterns of persecution” of
protective mothers that could then later be made part of a much larger
quantitative study of such phenomena. That is, by dissecting the family
court institution and its ravaging effects on mothers and children with
the kind of methodological scrutiny we employed, the often-opaque
configurations of a woefully malfunctioning arcane system were made
visible to us. Among those configurations and sketches that came to the
surface were the patterns of corruption that allow the weaponization of
family courts against unsuspecting mothers and children and, perhaps
most notably, against women of color. \pfl{}


Fourth, although we instantiated family court madness in our first
edition by drawing from a broad demographic mix, which naturally
included cases in the BIPOC population, in our new edition we specif-
ically analyze the effects of racism on the mistreatment of mothers by
judges and by child protective services. \pfl{}


For example, we provide a comparative analysis of how the same
child welfare agency and, moreover the same medical expert handling
the case for that agency, demonstrated a harsher attitude toward the
BIPOC mother, in which every statement she made or step she took
was assessed, alas, in her disfavor. She even felt intimidated by her own
attorney (who she later fired), fearing his reaction to her taking her four-
year-old daughter to the hospital emergency room after an incident of
sexual assault during an overnight visitation with the father. We formed
our analytical conclusions based on our examination of court records
\notePage{lxx}
and child protective service case files and by paying close attention to
the narrations of BIPOC mothers and, sometimes, their family mem-
bers as well. All in all, racial discrimination adds another layer of insult
for women of color who are already fighting a blatantly gender-biased
family court and child welfare system. \pfl{}


In looking at the racist elements in the cases we studied, we were sur-
prised to see a dearth of professional literature on how racism plays
an important role in protective mother cases. That is, racial factors in
protective mother cases have not served as an avid topic of study, al-
though it certainly deserves rigorous examination. Furthermore, even
though there are a few women of color who have participated in pro-
tective mother support groups, they are nonetheless underrepresented
in such groups. Thus, we use the opportunity in writing this second
edition to give proper attention to the racial factors that are essen-
tial in studying family courts’ mistreatment of women and children.
Basically, when a system descends to its lowest form of human degra-
dation, as evidenced by unrestrained racial bigotry, this must serve
as a clarion call that something is terribly awry in the system writ
large. \pfl{}


Fifth, although in our first edition we closely examined how bogus
mental health theories emerged in family law proceedings, this new vol-
ume, which employs a longitudinal approach to the study of protective
mother cases, gave us the opportunity to study how a mother’s day-to-
day existence is turned upside down by the quack mental health theories
that have come to haunt her. Mothers, engaged in recording, chron-
icling, and carefully documenting the sexual violations of their child,
have seen how the truth about those crimes can be so easily, almost in-
credulously, obliterated by the junk science diagnoses ascribed to them.
This is perhaps the greatest betrayal of battered women by the family
courts. \pfl{}


What’s even worse is that, as we show in this second edition, the
underpinnings of these new renditions of the archetypical “parental
alienation” theories are even more outlandish than their archetypes.
Analyzing the fundaments of Parental Alienation Syndrome (PAS) per-
mutations popularly circulating in the family courts today, we show
\notePage{lxxi}
how junk science theories aimed at discrediting protective mothers have
become increasingly more insidious over the years. \pfl{}


All things considered, we hope this second edition, by copiously doc-
umenting the institutional violence perpetrated against women in the
family courts, will inspire a major overhaul of our nation’s family courts.
This requires a concerted effort by the United States Department of Jus-
tice, the US Department of Health and Human Services (HHS), and the
US House of Representatives and Senate. Inasmuch as state legislatures
and state judicial conduct commissions have tried for four decades to
reform the courts without success, we sorely need the arm of the federal
government. The situation is dire: far too many mothers and children
are stuck in the quagmire of venal family courts, and they are paying
with their lives. \pfl{}


The second edition is organized as follows: We provide a detailed
supplement to the core chapters in the second section of the book:
Robed Rage, Lawless Law Guardians, Anti-Social Services, and Mental
Health Quackery. In addition, we provide supplementary material to the
third section of the book which focuses on reforms. All supplementary
material may be found under the subheading “Addendum.”


